\chapter{Создание датасета}
\label{ch:dataset}

Эта глава описывает путь от открытых игр на Ren'Py до таблицы пар ``текстовый эмбеддинг --- вектор музыки'', пригодной для обучения~\cite{nikolenko2018} и для расчёта метрик подстановки трека внутри одной игры.

\section{Извлечение текста и команд звука из сценария Ren'Py}

Сценарий игры представляет собой текстовые файлы со смесью диалогов вида \texttt{character ``реплика''}, описаний сцены и \textit{директив движка} --- команд воспроизведения и остановки музыки, переходов между треками, громкости и зацикливания. Такая структура удобна для автоматизации: можно последовательно читать файл, поддерживать текущее состояние ``какой трек считается активным'' и при появлении нового текстового блока записывать пару из текста и идентификатора трека.

Парсер в репозитории проекта~\appa{1} ориентирован на типичные конструкции Ren'Py~\cite{srcrenpydocs} (\texttt{play music}, \texttt{stop music}, вложенные метки и пр.). На практике разные игры слегка отличаются стилем написания сценария, поэтому для части проектов добавляются узкие исключения или постобработка. Корректность привязки ``абзац --- музыка'' проверяется выборочно; очевидные сбои (например, когда одна и та же дорожка висит на всём акте без обновления состояния парсера) отлавливаются по журналам прогона парсера и правилам фильтрации.

Точкой входа для пакетного разбора служит модуль \path{parse_rpy.py}: по каждой игре выполняются сканирование каталога (\texttt{scanGame}), восстановление связей по меткам (\texttt{traceGame}) и выгрузка CSV; при согласии оператора результат дополнительно загружается в PostgreSQL (см.\ раздел ниже про \texttt{insertDataset}).

\section{Парсер в коде: регулярные выражения и машина состояний}

Для устойчивого распознавания директив поверх разнородного текста без полноценного AST Ren'Py в проектной ветке используется набор регексов над строками \path{.rpy}. Например, команды включения дорожки на каналах \texttt{music}, \texttt{music2}, \texttt{bgs}, \texttt{sfx} сопоставляются шаблону \texttt{RE\_PLAY}: допускается как прямое имя файла в кавычках, так и идентификатор переменной с последующей индексацией по словарю (типичный приём скриптов Ren'Py, где дорожку выбирают из таблицы). Вспомогательная функция \texttt{playAudioRaw} извлекает строку пути из совпадения (листинг~\ref{lst:regex-play}).

\begin{lstlisting}[language=Python, caption={Фрагмент \texttt{trace\_labels.py}: директивы \texttt{play}/\texttt{queue} и извлечение необработанной ссылки на аудио.},label=lst:regex-play]
RE_PLAY = re.compile(
    r'^\s*(?:play|queue)\s+(music|sound|audio|sfx|music2|bgs)\s+'
    r'(?:\[\s*)?'
    r'(?:"([^"]+)"|(\w+))'
    r'(.*)',
    re.IGNORECASE,
)
RE_SUBSCRIPT_TAIL = re.compile(
    r'^\s*\[\s*["\']([^"\']+)["\']\s*\]',
)

def playAudioRaw(m_play):
    path_str = m_play.group(2)
    var_name = m_play.group(3)
    tail = m_play.group(4) or ""
    if path_str:
        return path_str
    if var_name:
        st = RE_SUBSCRIPT_TAIL.match(tail)
        if st:
            return var_name + '["' + st.group(1) + '"]'
        return var_name
    return ""
\end{lstlisting}

Полное движение по игре реализует \texttt{traceGame}: при обходе графа меток (\texttt{jump}, \texttt{return}, блоки перевода \texttt{translate}) поддерживается ассоциативное состояние каналов. В листинге~\ref{lst:parser-state} приведены ключевые поля состояния: фоновая музыка, зацикленность (\texttt{noloop} и \texttt{loop=False} разбираются отдельно), последняя «сырая» строка директивы (\texttt{raw\_music\_line}) для отладки и отложенная звуковая дорожка \texttt{SFX}.

\begin{lstlisting}[language=Python, caption={Фрагмент \texttt{trace\_labels.py}: состояние звуковых каналов в начале \texttt{traceGame}.},label=lst:parser-state]
NO_MUSIC = "NO_MUSIC"

state = {
    "music": NO_MUSIC,
    "music_loop": True,
    "raw_music_line": "",
    "bgs": NO_MUSIC,
    "bgs_loop": True,
    "raw_bgs_line": "",
    "pending_sfx": None,
    "sfx_invalidated": False,
}
\end{lstlisting}

При совпадении \texttt{RE\_PLAY} сначала резолвится путь через карту \texttt{audio\_map} и связанный словарь объявлений переменных, затем в зависимости от канала обновляется либо \texttt{music}/\texttt{bgs}, либо ставится очередное \texttt{pending\_sfx}. Параллельно обрабатываются переводческие блоки и вызовы \texttt{stop music}: это даёт возможность не терять смены саундтрека между языковыми вариантами сценария.

\section{Ручная разметка и источники шума}

Даже при успешном разборе скрипта остаются источники несоответствий между текстом и ``правильным'' для приложения музыкальным фоном:

\begin{itemize}
	\item автор игры может использовать один и тот же файл музыки для разных по настроению сцен --- модель видит только техническую связку, а не художественное намерение;
	\item границы текстовых блоков в данных могут не совпадать с интуитивными абзацами для читателя (слишком короткие или длинные отрезки);
	\item для части строк текст может быть техническим (меню, системные сообщения), их иногда исключают фильтрами.
\end{itemize}

Отдельный слой работы --- согласование идентификаторов треков между таблицей пар и таблицей признаков \texttt{music\_data}: файл после парсинга должен быть найден или сопоставлен с записью, для которой уже посчитан вектор $y$ (см.\ также обсуждение шума и утечки контекста при разметке~\cite{nikolenko2018}). Если аудио отсутствует или повреждено, строка исключается или заменяется.

\section{PostgreSQL как хранилище таблицы диалогов}

Чтобы поддерживать итерационную доработку правил парсера и переисчёт признаков без потери промежуточных результатов, пары текста и метаданные музыкального фона сохраняют в PostgreSQL (\texttt{psycopg}, строка подключения читается из переменной окружения \texttt{pghost}). Справочник треков \texttt{music\_data}(\texttt{id}, \texttt{game}, \texttt{path}) отделён от основной таблицы диалогов: в каждой строке сценария хранится внешний ключ \texttt{music} на единственный файл дорожки, что упрощает дедупликацию и многократный пересчёт акустики.

Основная таблица \texttt{text\_musics\_rel} (см.\ модуль \path{consts.py}) включает поля имени игры, файла, номера строки, спикера, оригинального и переведённого текста, признаков зацикленности музыки, сырой строки директивы и произвольного JSON в \texttt{extra}. Уникальный первичный ключ по тройке \texttt{(game, file, line\_number)} гарантирует адресование фиксированной реплики при повторных импортах CSV.

При разовой записи уже подготовленного списка словарей (последовательность строк датасета в памяти) использован потоковый \texttt{COPY FROM STDIN} через API курсора: это быстрее, чем сотни тысяч отдельных \texttt{INSERT} (листинг~\ref{lst:copy-insert}). Дополнительно в том же модуле реализованы \texttt{pullCsv}/\texttt{ensureTables}: CSV поднимают во временную таблицу \texttt{text\_musics\_rel\_temp}, после чего выполняют массовое \texttt{INSERT \ldots{} ON CONFLICT DO UPDATE}: при совпадающих ключах обновляют все полезные столбцы, оставляя стабильной идентичность строки в файле игры --- удобно при исправлении текста перевода без смены нумерации.

\begin{lstlisting}[language=Python, caption={Фрагмент \texttt{csv\_sync\_db.py}: пакетная загрузка датасета в PostgreSQL методом \texttt{copy}.},label=lst:copy-insert]
def insertDataset(dataset):
    ensureTables()
    conn = psycopg.connect(os.environ["pghost"])
    cur = conn.cursor()

    cols = list(dataset[0].keys())
    col_names = ", ".join(cols)
    sql = "copy " + TABLE_NAME + " (" + col_names + ") from stdin"

    with cur.copy(sql) as copy:
        for row in dataset:
            copy.write_row([row[col] for col in cols])

    conn.commit()
    conn.close()
\end{lstlisting}

Инструментальный запуск произвольного SQL доступен утилитой \path{music_scripts/db_exec.py}; служебные скрипты машинного обучения могут журналировать эксперимент и обновлять эмбеддинги в таблице \texttt{text\_musics\_rubert} уже после заполнения \texttt{text\_musics\_rel}.

\section{Формирование таблицы признаков и эмбеддингов текста}

После того как для каждой пригодной строки известны текст и \texttt{music\_id}, текст пропускается через модель RuBERT~\cite{kuratov2019} (семейство BERT~\cite{devlin2018}) или берётся из заранее посчитанного столбца в базе данных проекта. Так формируется входной вектор фиксированной размерности на строку.

\subsection{От моно-файла к широкой таблице признаков: иерархия процессоров}

\begin{sloppypar}
Целевая описательная строка дорожки $y\in\mathbb{R}^{d_y}$ составляется не одним огромным скриптом, а \textit{семейством независимых блоков признаков} в модуле \path{music_scripts/audioprocessors/}. Базовый класс \path{BaseAudioProcessor} задаёт неформально ``протокол'' расширения: каждый блок обязан перечислить целевые столбцы с типами PostgreSQL (\path{DBColumnDes}) и реализовать метод \path{process}, принимающий сигнал $y$, частоту дискретизации~\path{sr}, путь к файлу и контекст. Дополнительно класс задаёт офлайновую нормализацию z-score по таблице; перечень колонок и режим задаёт метод \path{column_normalization_specs}~(листинг~\ref{lst:base-ap}); для столбцов с тяжёлыми хвостами допускается предобразование $\ln(1+x)$ перед вычитанием среднего\@.
\end{sloppypar}

\begin{lstlisting}[language=Python, caption={Фрагмент \texttt{BaseAudioProcessor.py}: контракт процессора и объявление схемы столбца.},label=lst:base-ap]
class BaseAudioProcessor:
    @property
    def processor_key(self) -> str:
        return self.__class__.__name__

    @property
    def target_columns(self) -> list[DBColumnDes]:
        raise NotImplementedError

    def process(self, y: np.ndarray, sr: int, path: str, ctx: dict) -> dict:
        raise NotImplementedError

    def column_normalization_specs(self) -> tuple[ColumnNormalizationSpec, ...]:
        return ()
\end{lstlisting}

\begin{sloppypar}
Конкретные реализации (\texttt{MFCC\-Processor}, \texttt{Spec\-tral\-Processor}, \texttt{Basic\-Stats\-Processor} и~др.)\ добавляются в список \texttt{processors} объекта~\texttt{Audio\-En\-rich\-er}: по файлу они считывают волну через~\texttt{librosa}~\cite{librosa}, цепочка процессоров затем переводится в столбцы таблицы операцией~\texttt{UPDATE} языка~SQL\@. Листинг~\ref{lst:mfcc} иллюстрирует схему: после перевода мел-спектрограммы в деци\-бель\-ную шкалу коэффициенты MFCC покад\-ро\-во агрегируются по сред\-нему, ста\-нд\-арт\-ному от\-клоне\-ни\-ю и кван\-ти\-лям уровней $5$\% и~$95$\% для каж\-до\-го ко\-эффи\-ци\-ен\-та; тем самым задаётся компактный набор статистик почти посто\-ян\-ной раз\-мер\-но\-сти без жёст\-кой связи с абсо\-лют\-ной дли\-те\-ль\-нос\-тью дорож\-ки\@.
\end{sloppypar}

\begin{lstlisting}[language=Python, caption={Фрагмент \texttt{MFCCProcessor.py}: признаки MFCC после мел-представления.},label=lst:mfcc]
class MFCCProcessor(BaseAudioProcessor):
    def process(self, y: np.ndarray, sr: int, path: str, ctx: dict) -> dict:
        S = librosa.feature.melspectrogram(
            y=y, sr=sr, n_fft=N_FFT, hop_length=HOP_LENGTH,
            n_mels=N_MELS, power=2.0)
        S_db = librosa.power_to_db(S)
        mfcc = librosa.feature.mfcc(S=S_db, n_mfcc=N_MFCC)
        mfcc_mean = np.mean(mfcc, axis=1)
        mfcc_std = np.std(mfcc, axis=1)
        mfcc_p05 = np.percentile(mfcc, 5, axis=1)
        mfcc_p95 = np.percentile(mfcc, 95, axis=1)
        out = {}
        for i in range(N_MFCC):
            out[f"mfcc_mean_{i:02d}"] = float(mfcc_mean[i])
            out[f"mfcc_std_{i:02d}"] = float(mfcc_std[i])
            out[f"mfcc_p05_{i:02d}"] = float(mfcc_p05[i])
            out[f"mfcc_p95_{i:02d}"] = float(mfcc_p95[i])
        return out
\end{lstlisting}

Результаты сохраняются в столбцах строки \texttt{music\_data} для каждого трека, что упрощает контроль: достаточно пересобрать отдельные процессоры, не трогая таблицу диалогов, пока ключи дорожек стабильны.

Итог --- большая таблица строк вида: игра, порядковый контекст внутри игры, эмбеддинг текста, идентификатор музыки, ссылка на вектор $y$. В типичном сохранённом прогоне насчитывается порядка $1{,}9\cdot 10^4$ строк при $d_y = 320$.

\section{Каталог кандидатов для проверки подстановки трека}

Метрика ``попадание в топ доли кандидатов'', используемая в главе~\ref{ch:training}, требует для каждой тестовой строки знать все треки, которые в принципе встречаются в данной игре, и их векторы $y$. Этот каталог может храниться отдельно на диске (например, в каталоге \texttt{candidates/} в составе экспорта для PyTorch), либо вычисляться запросом к PostgreSQL через объединение \texttt{text\_musics\_rel} с \texttt{music\_data} при соблюдении тех же ограничений по игре. Важно, чтобы множество кандидатов совпадало с тем, что видит пользователь при выборе из саундтрека игры; иначе оценка будет завышена или занижена.

\section{Хранение и воспроизводимость}

\begin{sloppypar}
Для каждого эксперимента фиксируются версия таблицы эмбеддингов, состав столбцов $y$, список игр и параметры фильтрации строк. Это позволяет сопоставлять численные результаты между архитектурами; конкретные идентификаторы прогонов и имена файлов [\path{config.json},~\path{metrics.json}] указаны в манифесте \path{manifest.jsonl}~\appa{3} репозитория~\appa{1}.
\end{sloppypar}
